\section{Introduction}

Learning compact dynamical models is central to scientific machine learning when
high-fidelity simulations are too expensive to evaluate repeatedly.
The challenge becomes substantially harder when a physical parameter changes not only the
magnitude of the response, but also the qualitative character of the dynamics.
Parameterized incompressible flows provide a representative example: different portions of
the parameter domain may exhibit relaxation toward a steady state, weak oscillatory growth
near instability onset, or sustained nonlinear periodic motion.
A reduced model intended for many-query prediction must therefore remain accurate under
autonomous rollout while the underlying dynamics vary across qualitatively different
regimes.

This setting exposes two coupled difficulties.
First, the state representation itself can be strongly regime dependent.
A low-dimensional basis that efficiently describes nearly steady states need not provide an
equally compact description of developed periodic dynamics.
Second, even when a suitable reduced representation is available, the corresponding
dynamical model must reproduce different temporal mechanisms, including decay, growth,
phase evolution, and nonlinear saturation.
A single global reduced model must solve both problems simultaneously: it must encode
states drawn from distinct local geometries and learn a dynamical law that remains accurate
over all of them.
This requirement is particularly demanding for autonomous prediction, where small local
errors are recursively propagated over long horizons.

Projection-based reduced-order models provide an explicit low-dimensional dynamical
structure through methods such as proper orthogonal decomposition (POD) and Galerkin
projection \citep{sirovich1987,berkooz1993,benner2015}.
Their predictive accuracy, however, can deteriorate when truncated modes contribute to
dissipation, nonlinear interactions, or other unresolved effects \citep{amsallem2012}.
Data-driven reduced models instead learn the dynamics from trajectories and can represent
richer nonlinear mappings \citep{hesthaven2018,pathak2018,vlachas2018,lu2021,kovachki2023}.
Yet accurate one-step prediction does not in general imply a stable recursive dynamical
model, because prediction errors modify the inputs supplied to subsequent steps
\citep{duraisamy2019,brunton2020}.
Hybrid approaches combine projected physical structure with learned corrections and provide
a natural way to retain known dynamics while modeling unresolved components
\citep{dar2023,lee2020,fresca2022,geelen2023}.

A complementary strategy is to localize the reduced representation.
Local subspaces, clustered reduced models, and nonlinear local manifolds can specialize to
different portions of a parameterized solution set \citep{kaiser2014,li2021cluster,diaz2024}.
Localization, however, introduces a difficulty that is less visible in a global model.
Two independently constructed local reduced models generally have different affine means,
bases, scalings, and reduced operators.
Their reduced coordinates therefore do not possess a canonical componentwise
correspondence.
Consequently, switching or averaging local coefficient vectors can mix incompatible
representations.
This raises a broader learning problem: how should multiple local dynamical models be
selected and combined when each model evolves in its own reduced coordinate system?

Mixture-of-experts (MoE) models offer a natural language for conditional specialization.
Classical and modern MoE architectures route inputs to subsets of specialized predictors,
enabling conditional computation and expert specialization
\citep{jacobs1991,jordan1994,shazeer2017,fedus2022}.
For regime-dependent dynamical systems, however, two different forms of specialization are
needed.
Within a fixed regime, unresolved dynamics may benefit from multiple conditional closure
experts.
Across regimes, complete reduced dynamical models must themselves be selected or combined.
The latter problem differs from a conventional MoE layer because the candidate specialists
need not share a common reduced state representation.

We introduce the \emph{Regime-Adaptive Localized Mixture-of-Experts Reduced-Order Model}
(RAL-MoE-ROM), a two-level hierarchical MoE architecture for this setting.
At the representation level, the parameterized solution set is covered by overlapping
regime-local reduced charts.
Each chart supports a complete autonomous physics--data specialist whose projected reduced
dynamics are augmented by a structured sparse MoE closure.
The internal MoE uses hierarchical sparse routing to select among experts with nonlinear,
affine, and low-rank bilinear branches, allowing the learned correction to adapt within a
local dynamical regime while preserving the projected model as an explicit inductive bias.

A second routing level operates across the regime-local specialists.
The outer router determines whether a single specialist or an adjacent specialist pair
should contribute to the prediction.
Crucially, the selected specialists are not mapped into a common latent coordinate system.
They are initialized from the same physical history, encoded independently in their native
charts, and advanced independently under their own autonomous reduced dynamics.
When two adjacent specialists are used, the proposed Top-2 convex correction (T2-C)
combines their reconstructed physical fields rather than their modal states or reduced
right-hand sides.
This separation between within-specialist expert routing and between-specialist model
routing is illustrated in Fig.~\ref{fig:ral_overview}.

Our contributions are threefold:
\begin{itemize}
    \item \textbf{Two-level hierarchical expert specialization.}
    We formulate regime-adaptive reduced dynamics as a hierarchy in which sparse expert
    routing models unresolved dynamics within each local specialist, while a second router
    selects or combines complete specialists across dynamical regimes.

    \item \textbf{Structured sparse MoE closures for autonomous reduced dynamics.}
    Each specialist augments projected dynamics with conditionally activated expert maps
    containing nonlinear, affine, and low-rank bilinear components.
    The resulting model preserves an explicit physics-based backbone while adapting the
    learned closure to the local reduced trajectory.

    \item \textbf{Physical-space assembly of noncommensurate local models.}
    Adjacent specialists retain independent reduced coordinates and autonomous rollouts.
    Cross-regime coupling is performed only after physical reconstruction through T2-C,
    avoiding direct interpolation of incompatible modal states or reduced operators.
\end{itemize}

\begin{figure}[t]
    \centering
    \resizebox{\linewidth}{!}{\begin{tikzpicture}[
  font=\scriptsize,
  node distance=4mm and 5mm,
  block/.style={draw,rounded corners=1.5pt,align=center,minimum height=7mm,inner sep=3pt,fill=blue!4},
  chart/.style={draw,rounded corners=2pt,align=center,minimum width=46mm,minimum height=14mm,inner sep=3pt},
  route/.style={draw,rounded corners=1.5pt,align=center,minimum height=7mm,inner sep=3pt,fill=orange!12},
  outputblock/.style={draw,rounded corners=1.5pt,align=center,minimum height=7mm,inner sep=3pt,fill=green!10},
  arrow/.style={-{Latex[length=1.8mm]},thick},
  soft/.style={-{Latex[length=1.5mm]},thin,dashed}
]
\node[block] (history) {physical history\\$\{\bm u,\bm p\}$ and $\mu$};
\node[route,above=7mm of history] (outer) {Level 2\\regime router};

\node[chart,right=12mm of history] (hopf) {\textbf{Hopf chart}: native encoding $\rightarrow$ projected backbone\\Level-1 sparse experts $\rightarrow$ autonomous rollout $\rightarrow$ decode};
\node[chart,above=3mm of hopf] (steady) {\textbf{steady chart}: native encoding $\rightarrow$ projected backbone\\Level-1 sparse experts $\rightarrow$ autonomous rollout $\rightarrow$ decode};
\node[chart,below=3mm of hopf] (periodic) {\textbf{periodic chart}: native encoding $\rightarrow$ projected backbone\\Level-1 sparse experts $\rightarrow$ autonomous rollout $\rightarrow$ decode};

\node[outputblock,right=10mm of hopf] (fusion) {Top-1 or T2-C\\physical-space fusion};
\node[outputblock,right=5mm of fusion] (prediction) {predicted\\physical fields};

\draw[arrow] (history) -- (outer);
\draw[arrow] (outer.east) -- (steady.west);
\draw[arrow] (outer.east) -- (hopf.west);
\draw[arrow] (outer.east) -- (periodic.west);
\draw[arrow] (steady.east) -- (fusion.west);
\draw[arrow] (hopf.east) -- (fusion.west);
\draw[arrow] (periodic.east) -- (fusion.west);
\draw[arrow] (fusion) -- (prediction);

\node[above=1mm of steady.north,align=center] {independent native coordinates};
\end{tikzpicture}
}
    \caption{Overview of RAL-MoE-ROM. The method uses sparse expert routing within each regime-local specialist and regime-level routing across complete local dynamical models. Adjacent specialists evolve independently in their native reduced coordinates and are combined only after reconstruction in the common physical space.}
    \label{fig:ral_overview}
\end{figure}

We evaluate RAL-MoE-ROM on three parametrized incompressible-flow systems spanning steady,
Hopf-onset, and periodic dynamics.
The experiments are designed to separate the roles of structured expert architecture,
projected physical dynamics, regime localization, and cross-regime routing.
Across the evaluated long-horizon rollouts, the local specialists remain accurate and
stable, while T2-C consistently improves the corresponding hard or single-specialist
prediction in the tested overlap regions.
The results also show that localization is not uniformly superior to a global model at every
parameter value; rather, its value emerges from combining local representation,
regime-specific dynamics, and adaptive cross-regime assembly.
